\chapter{Algebraic Groups and Abelian Varieties}\label{algebraic-groups}
\chapterstatus{complete}

\section{Introduction}\label{algebraic-groups:section-introduction}

A Hodge structure on a real vector space $V$ is the same as an action of the real algebraic group $\mathbb{S} = \mathrm{Res}_{\CC/\RR}\mathbb{G}_m$, the
Deligne torus (Proposition~\ref{algebraic-groups:proposition-deligne-torus}). The symmetry group of a rational Hodge structure, its Mumford--Tate group
(Chapter~\ref{hodge-structures}), is the smallest algebraic subgroup of $\mathrm{GL}(V)$ defined over $\QQ$ through which this action factors. This chapter provides
the language: affine group schemes and their Hopf algebras, the embedding of every algebraic group into some $\mathrm{GL}_n$, representations, tori and characters,
Lie algebras, reductive groups and Weil restriction. The ground field is a field $k$ of characteristic $0$ unless stated otherwise. References:
\cite{waterhouse-affine}, \cite{milne-algebraic-groups}, \cite{borel-lag}, \cite{deligne-milne}.

\noindent\textbf{Prerequisites.} Chapter~\ref{schemes} (Schemes), Chapter~\ref{lie-groups} (Lie Groups) and Chapter~\ref{representation-theory} (Representations of
Groups and Lie Algebras).

\section{Affine group schemes}\label{algebraic-groups:section-group-schemes}

\begin{definition}\label{algebraic-groups:definition-group-scheme}
An \emph{affine group scheme} over $k$ is an affine $k$-scheme $G = \Spec A$ whose functor of points $R \mapsto G(R) = \Hom_{k\text{-alg}}(A, R)$ takes values in groups,
functorially in the $k$-algebra $R$. Equivalently (Yoneda), $A = \mathcal{O}(G)$ is a \emph{commutative Hopf algebra}: it has $k$-algebra maps
$\Delta \colon A \to A \otimes A$ (comultiplication), $\varepsilon \colon A \to k$ (counit) and $S \colon A \to A$ (antipode), dual to multiplication, identity and inverse,
satisfying the dual group axioms. A \emph{linear algebraic group} is an affine group scheme of finite type. Homomorphisms are morphisms of schemes compatible with the
group laws.
\end{definition}

\begin{example}\label{algebraic-groups:example-groups}
\begin{enumerate}
\item The \emph{additive group} $\mathbb{G}_a = \Spec k[t]$, $\Delta(t) = t \otimes 1 + 1 \otimes t$, with $\mathbb{G}_a(R) = (R, +)$.
\item The \emph{multiplicative group} $\mathbb{G}_m = \Spec k[t, t^{-1}]$, $\Delta(t) = t \otimes t$, with $\mathbb{G}_m(R) = (R^*, \cdot)$. Its subgroup $\mu_n = \Spec k[t]/(t^n - 1)$ of
$n$-th roots of unity.
\item $\mathrm{GL}_n = \Spec k[x_{ij}, 1/\det]$ with $\Delta(x_{ij}) = \sum_lx_{il} \otimes x_{lj}$, and its closed subgroups $\mathrm{SL}_n$ ($\det = 1$), the orthogonal group
$\mathrm{O}(q) = \{g : q(gv) = q(v)\}$ of a nondegenerate quadratic form, the symplectic group $\mathrm{Sp}_{2g}$ of a nondegenerate alternating form, and the
diagonal torus $\mathbb{G}_m^n$.
\item For a finite-dimensional vector space $V$, $\mathrm{GL}(V)(R) = \mathrm{Aut}_R(V \otimes_kR)$.
\end{enumerate}
\end{example}

\begin{definition}\label{algebraic-groups:definition-representation}
A \emph{representation} of $G$ on a $k$-vector space $V$ is a natural action of $G(R)$ on $V \otimes_kR$ by $R$-linear maps. Equivalently it is a \emph{comodule}: a linear map
$\rho \colon V \to V \otimes A$ with $(\rho \otimes \mathrm{id})\rho = (\mathrm{id} \otimes \Delta)\rho$ and $(\mathrm{id} \otimes \varepsilon)\rho = \mathrm{id}$. The element $g \in G(R)$,
that is $g \colon A \to R$, acts by $(\mathrm{id} \otimes g) \circ \rho$ extended $R$-linearly. The \emph{regular representation} is $A$ itself with $\rho = \Delta$.
\end{definition}

\begin{lemma}\label{algebraic-groups:lemma-locally-finite}
Every comodule is the union of its finite-dimensional subcomodules.
\end{lemma}

\begin{proof}
Let $v \in V$ and choose a basis $(a_i)$ of $A$. Write $\rho(v) = \sum_iv_i \otimes a_i$ (finite sum) and $\Delta(a_i) = \sum_{j,l}c_{ijl}a_j \otimes a_l$. Applying the two sides of
the coassociativity identity to $v$ gives $\sum_i\rho(v_i) \otimes a_i = \sum_{i,j,l}c_{ijl}v_i \otimes a_j \otimes a_l$, and comparing the coefficients of $a_l$ in the last factor gives
$\rho(v_l) = \sum_{i,j}c_{ijl}v_i \otimes a_j$. So $W = \mathrm{span}(v_i)$ is a subcomodule, and $v = \sum_i\varepsilon(a_i)v_i \in W$ by the counit identity.
\end{proof}

\begin{theorem}\label{algebraic-groups:theorem-linear}
Every linear algebraic group $G$ is isomorphic to a closed subgroup of some $\mathrm{GL}_n$.
\end{theorem}

\begin{proof}
Let $A = \mathcal{O}(G)$ with algebra generators $f_1, \ldots, f_m$. By Lemma~\ref{algebraic-groups:lemma-locally-finite} for the regular representation, there is a
finite-dimensional subcomodule $V \subseteq A$ containing all $f_i$. Choose a basis $v_1, \ldots, v_n$ of $V$ and write $\Delta(v_j) = \sum_iv_i \otimes a_{ij}$ with $a_{ij} \in A$.
Coassociativity shows that $x_{ij} \mapsto a_{ij}$ defines a Hopf algebra map $k[x_{ij}, 1/\det] \to A$, that is, a homomorphism $\rho \colon G \to \mathrm{GL}_n$ (the determinant maps
to a unit, whose inverse is the image of the determinant of $S(a_{ij})$). It is surjective on rings, because $v_j = (\varepsilon \otimes \mathrm{id})\Delta(v_j) = \sum_i\varepsilon(v_i)a_{ij}$ lies in
the image, and the $v_j$ include the generators $f_i$. So $\rho$ is a closed immersion.
\end{proof}

\begin{remark}\label{algebraic-groups:remark-smoothness}
In characteristic $0$ every linear algebraic group is smooth (Cartier's theorem) \cite{milne-algebraic-groups}, so its reduced structure causes no trouble and
$\dim G = \dim\mathrm{Lie}(G)$. In characteristic $p$, $\mu_p = \Spec k[t]/(t^p - 1) = \Spec k[t]/(t - 1)^p$ is a non-reduced group scheme with one point.
\end{remark}

\section{Representations and Tannakian categories}\label{algebraic-groups:section-tannakian}

\begin{definition}\label{algebraic-groups:definition-tensor-category}
The finite-dimensional representations of $G$ form an abelian category $\mathrm{Rep}_k(G)$ with tensor products, duals and the unit object $k$ (trivial representation), and
a \emph{fibre functor} $\omega \colon \mathrm{Rep}_k(G) \to \mathrm{Vec}_k$ forgetting the action. For a representation $V$, the tensor spaces are
$T^{r,s}(V) = V^{\otimes r} \otimes (V^\vee)^{\otimes s}$, and a \emph{tensor} $t \in T^{r,s}(V)$ is \emph{fixed} by $G$ if $g(t \otimes 1) = t \otimes 1$ for all $g \in G(R)$ and all $R$.
\end{definition}

\begin{theorem}\label{algebraic-groups:theorem-tannakian}
\begin{enumerate}
\item (Tannaka duality) $G$ is recovered from $(\mathrm{Rep}_k(G), \omega)$: $G(R)$ is the group of families $(\lambda_V)$ of $R$-linear automorphisms of $V \otimes R$,
natural in $V$ and compatible with tensor products and the unit.
\item (Chevalley) Every closed subgroup $H \subseteq G$ is the stabiliser of a line $L$ in some finite-dimensional representation $W$ of $G$.
\item If $H \subseteq \mathrm{GL}(V)$ is reductive (Definition~\ref{algebraic-groups:definition-reductive}), then $H$ is the subgroup of $\mathrm{GL}(V)$ fixing all tensors in
$\bigoplus_{r,s}T^{r,s}(V)$ that are fixed by $H$.
\end{enumerate}
\end{theorem}

Statement (1) is \cite[Theorem 2.11]{deligne-milne}, (2) is Chevalley's theorem \cite{milne-algebraic-groups}, and (3) is
\cite[Proposition 3.1(c)]{deligne-abelian}. Statement (3) is the reason why Mumford--Tate groups, defined as stabilisers of Hodge tensors, capture all the Hodge classes in
tensor spaces (Chapter~\ref{hodge-structures}). Without reductivity (3) fails. Every tensor fixed by the group $B$ of invertible upper triangular $2 \times 2$ matrices is fixed by
$\mathrm{GL}_2$: a $B$-fixed vector in a representation of $\mathrm{GL}_2$ is a highest weight vector of weight $0$, so it spans a trivial subrepresentation. But $B \neq \mathrm{GL}_2$.

\begin{proposition}\label{algebraic-groups:proposition-subgroup-fixing}
Let $G$ be a linear algebraic group and $V$ a representation. For $t \in T^{r,s}(V)$ the stabiliser $G_t = \{g : gt = t\}$ is a closed subgroup of $G$, defined over the
field of definition of $t$. For $v \in V$, the stabiliser of the line $kv$ is a closed subgroup.
\end{proposition}

\begin{proof}
Write $\rho(t) = \sum_it_i \otimes a_i$ with the $t_i$ a basis of $T^{r,s}(V)$ containing $t = t_1$. Then $g \in G(R)$ fixes $t$ if and only if $g(a_1) = 1$ and $g(a_i) = 0$ for $i \geq 2$,
which are polynomial equations. For the line, $g$ stabilises $kv$ if and only if $gv \wedge v = 0$ in $\Lambda^2V$, again polynomial equations.
\end{proof}

\section{Tori, characters and the Deligne torus}\label{algebraic-groups:section-tori}

\begin{definition}\label{algebraic-groups:definition-torus}
A \emph{character} of $G$ is a homomorphism $G \to \mathbb{G}_m$; they form an abelian group $X^*(G)$ under pointwise multiplication (written additively). A
\emph{cocharacter} is a homomorphism $\mathbb{G}_m \to G$. A \emph{split torus} is a group isomorphic to $\mathbb{G}_m^r$, and a \emph{torus} is a group $T$ such that $T_{\bar k}$ is a split
torus. $T$ is \emph{split} over a field extension $L$ if $T_L$ is.
\end{definition}

\begin{proposition}\label{algebraic-groups:proposition-torus-reps}
Let $T = \mathbb{G}_m^r$. Then $X^*(T) \cong \ZZ^r$, the characters being the monomials $t^m = t_1^{m_1} \cdots t_r^{m_r}$, and $\mathcal{O}(T) = k[X^*(T)]$ is the group algebra. A
representation of $T$ is the same as a $X^*(T)$-graded vector space $V = \bigoplus_\chi V_\chi$, where $T$ acts on $V_\chi$ through $\chi$. In particular every representation of a
split torus is a direct sum of one-dimensional ones.
\end{proposition}

\begin{proof}
$\mathcal{O}(T) = k[t_i^{\pm 1}]$ has basis the monomials $t^m$, which are group-like: $\Delta(t^m) = t^m \otimes t^m$. A character is a group-like unit, and the group-like elements of a
group algebra are the group elements (a group-like element $\sum c_mt^m$ satisfies $\sum c_mt^m \otimes t^m = \sum c_mc_{m'}t^m \otimes t^{m'}$, forcing a single term with $c_m = 1$). For a
comodule $V$ put $V_m = \{v : \rho(v) = v \otimes t^m\}$. Writing $\rho(v) = \sum_mv_m \otimes t^m$, coassociativity gives $\rho(v_m) = v_m \otimes t^m$, so $v_m \in V_m$, and the counit gives
$v = \sum_mv_m$. The sum of the $V_m$ is direct because the $t^m$ are linearly independent.
\end{proof}

\begin{remark}\label{algebraic-groups:remark-galois-tori}
For a torus $T$ over $k$, split over a finite Galois extension $L$, $\mathrm{Gal}(L/k)$ acts on $X^*(T_L) \cong \ZZ^r$, and $T \mapsto X^*(T_L)$ is an anti-equivalence between tori
over $k$ split by $L$ and free abelian groups of finite rank with a $\mathrm{Gal}(L/k)$-action \cite{milne-algebraic-groups}. Representations of $T$ are
$X^*(T_L)$-graded $L$-vector spaces $V_L$ with a semilinear Galois action compatible with the grading.
\end{remark}

\begin{definition}\label{algebraic-groups:definition-deligne-torus}
The \emph{Deligne torus} is $\mathbb{S} = \mathrm{Res}_{\CC/\RR}\mathbb{G}_m$ (Section~\ref{algebraic-groups:section-weil-restriction}), the real algebraic group with
$\mathbb{S}(R) = (R \otimes_\RR\CC)^*$. Concretely $\mathbb{S} = \Spec\RR[a, b, 1/(a^2 + b^2)]$, $\mathbb{S}(\RR) = \CC^*$ via $z = a + ib$, and
$\mathbb{S}(\CC) = (\CC \otimes_\RR\CC)^* \cong \CC^* \times \CC^*$, where $z \in \CC^* = \mathbb{S}(\RR)$ maps to $(z, \bar z)$. The characters of $\mathbb{S}_\CC$ are
$z^p\bar z^q$, $(p, q) \in \ZZ^2$. The \emph{weight cocharacter} $w \colon \mathbb{G}_m \to \mathbb{S}$ is the inclusion $\RR^* \subseteq \CC^*$, and the \emph{norm}
$\mathrm{Nm} \colon \mathbb{S} \to \mathbb{G}_m$ is $z \mapsto z\bar z$.
\end{definition}

\begin{proposition}\label{algebraic-groups:proposition-deligne-torus}
Let $V$ be a finite-dimensional real vector space. Representations $h \colon \mathbb{S} \to \mathrm{GL}(V)$ correspond bijectively to decompositions
$V_\CC = \bigoplus_{p,q}V^{p,q}$ with $\overline{V^{p,q}} = V^{q,p}$, via
\[
V^{p,q} = \{v \in V_\CC : h(z)v = z^{-p}\bar z^{-q}v \text{ for all } z \in \CC^*\}.
\]
Under this correspondence $w(r)$ acts on $V^{p,q}$ by $r^{-(p+q)}$, so $V = \bigoplus_nV_n$ with $V_{n,\CC} = \bigoplus_{p+q=n}V^{p,q}$ is the \emph{weight decomposition}.
\end{proposition}

\begin{proof}
$\mathbb{S}_\CC$ is a split torus with character group $\ZZ^2$ generated by $z$ and $\bar z$. By Proposition~\ref{algebraic-groups:proposition-torus-reps}, a representation of $\mathbb{S}_\CC$ on
$V_\CC$ is a $\ZZ^2$-grading, which we index by $(p, q) \mapsto$ the character $z^{-p}\bar z^{-q}$ (Deligne's sign convention). The subgroup $\mathbb{S}(\RR) = \CC^*$ is Zariski dense in
$\mathbb{S}_\CC$, since a Laurent polynomial in $z$, $\bar z$ vanishing on $\{(z, \bar z)\}$ is zero, so the grading is determined by the action of $\CC^*$ as displayed. The
representation is defined over $\RR$ if and only if $V_\CC$ is stable under the action of $\CC^*$ and $h(z)$ commutes with complex conjugation on $V_\CC$. If $v \in V^{p,q}$, then
$h(z)\bar v = \overline{h(z)v} = \overline{z^{-p}\bar z^{-q}v} = z^{-q}\bar z^{-p}\bar v$, so $\bar v \in V^{q,p}$. Conversely, a grading with $\overline{V^{p,q}} = V^{q,p}$ defines $h(z)$
by the formula, and $h(z)$ commutes with conjugation by the same computation, hence is real. The statement about $w$ is the case $z = r \in \RR^*$.
\end{proof}

\section{The Lie algebra}\label{algebraic-groups:section-lie-algebra}

\begin{definition}\label{algebraic-groups:definition-lie-algebra}
Let $k[\varepsilon] = k[t]/(t^2)$. The \emph{Lie algebra} of $G$ is $\mathrm{Lie}(G) = \ker(G(k[\varepsilon]) \to G(k))$, a $k$-vector space identified with the tangent space
$(\mathfrak{m}_e/\mathfrak{m}_e^2)^\vee$ at the identity (Exercise~\ref{schemes:exercise-dual-numbers}). The bracket is induced by the commutator in
$G(k[\varepsilon_1, \varepsilon_2])$, and a representation $\rho$ of $G$ differentiates to a representation $d\rho$ of $\mathrm{Lie}(G)$.
\end{definition}

\begin{example}\label{algebraic-groups:example-lie-algebras}
$\mathrm{Lie}(\mathrm{GL}_n) = \{1 + \varepsilon X\} \cong \mathfrak{gl}_n$, with the commutator bracket. Since $\det(1 + \varepsilon X) = 1 + \varepsilon\operatorname{tr}X$,
$\mathrm{Lie}(\mathrm{SL}_n) = \{X : \operatorname{tr}X = 0\}$. Similarly $\mathrm{Lie}(\mathrm{O}(q)) = \{X : B(Xv, w) + B(v, Xw) = 0\}$, where $B$ is the bilinear form of $q$, and
$\mathrm{Lie}(\mathbb{G}_m^r) = k^r$ with zero bracket. For $k = \RR$ or $\CC$, $\mathrm{Lie}(G)$ is the Lie algebra of the Lie group $G(k)$
(Definition~\ref{lie-groups:definition-lie-algebra-of-g}).
\end{example}

\begin{proposition}\label{algebraic-groups:proposition-lie-correspondence}
In characteristic $0$, for connected $G$, $H \mapsto \mathrm{Lie}(H)$ is an injective, inclusion-preserving map from connected closed subgroups of $G$ to Lie subalgebras of
$\mathrm{Lie}(G)$. A subspace $W \subseteq V$ of a representation is $G$-stable if and only if it is $\mathrm{Lie}(G)$-stable, and a vector is fixed by $G$ if and only if it is
killed by $\mathrm{Lie}(G)$.
\end{proposition}

See \cite{borel-lag} and \cite{milne-algebraic-groups}. It fails in characteristic $p$: $\mu_p$ and $\{1\}$ have
different Lie algebras over a field of characteristic $p$ but the same points, and $\mathbb{G}_a$ acting on $k^2$ by $t \mapsto \left(\begin{smallmatrix}1 & t^p \\ 0 & 1\end{smallmatrix}\right)$
has $\mathrm{Lie}$ acting trivially without fixing everything.

\section{Reductive groups}\label{algebraic-groups:section-reductive}

\begin{definition}\label{algebraic-groups:definition-reductive}
A linear algebraic group $G$ over $k$ is \emph{unipotent} if it is isomorphic to a closed subgroup of the upper triangular unipotent matrices. The \emph{unipotent radical}
$R_u(G)$ is the largest connected normal unipotent subgroup of $G_{\bar k}$. A connected $G$ is \emph{reductive} if $R_u(G) = 1$, and \emph{semisimple} if moreover its centre is finite.
\end{definition}

\begin{theorem}\label{algebraic-groups:theorem-reductive}
Let $k$ have characteristic $0$ and $G$ be connected.
\begin{enumerate}
\item $G$ is reductive if and only if every finite-dimensional representation of $G$ is semisimple (a direct sum of irreducible representations).
\item If $G$ is reductive, then $G = Z(G)^\circ \cdot G^{\mathrm{der}}$, where the connected centre $Z(G)^\circ$ is a torus and the derived group $G^{\mathrm{der}}$ is semisimple, and
$Z(G)^\circ \cap G^{\mathrm{der}}$ is finite.
\item Tori, $\mathrm{GL}_n$, $\mathrm{SL}_n$, $\mathrm{Sp}_{2g}$ and $\mathrm{SO}(q)$ are reductive; $\mathbb{G}_a$ and the group of invertible upper triangular matrices are not.
\end{enumerate}
\end{theorem}

See \cite{milne-algebraic-groups} and \cite{borel-lag}. For example, the representation of $\mathbb{G}_a$ on $k^2$ by
$t \mapsto \left(\begin{smallmatrix}1 & t \\ 0 & 1\end{smallmatrix}\right)$ has the invariant line $k e_1$ without an invariant complement, so $\mathbb{G}_a$ is not reductive.

\begin{proposition}[Compact real forms]\label{algebraic-groups:proposition-cartan-involution}
Let $G$ be a connected linear algebraic group over $\RR$ and $\theta$ an involution of $G_\CC$ such that
$G^{(\theta)} = \{g \in G(\CC) : \theta(\bar g) = g\}$ is compact. Then $G$ is reductive. In particular, if $G \subseteq \mathrm{GL}(V)$ and there is a positive definite inner
product on $V$ such that $G$ is stable under $g \mapsto g^T$ (the adjoint), then $G$ is reductive.
\end{proposition}

\begin{proof}
The group $K = G^{(\theta)}$ is a compact real form of $G(\CC)$: it is the fixed group of the antiholomorphic involution $g \mapsto \theta(\bar g)$, so its Lie algebra $\mathfrak{k}$ is a real
form of $\mathrm{Lie}(G_\CC)$, that is, $\mathrm{Lie}(G_\CC) = \mathfrak{k} \oplus i\mathfrak{k}$. Let $W$ be a complex representation of $G_\CC$ and $W' \subseteq W$ a subrepresentation. By
Corollary~\ref{lie-groups:corollary-averaging}, $W$ has a $K$-invariant Hermitian inner product, and the orthogonal complement $W''$ of $W'$ is $K$-stable, hence $\mathfrak{k}$-stable. As
$\mathfrak{k}$ spans $\mathrm{Lie}(G_\CC)$ over $\CC$, $W''$ is $\mathrm{Lie}(G_\CC)$-stable, hence $G_\CC$-stable by Proposition~\ref{algebraic-groups:proposition-lie-correspondence}. So
every representation of $G_\CC$, and therefore of $G$, is semisimple, and $G$ is reductive by Theorem~\ref{algebraic-groups:theorem-reductive}(1). For the second statement,
extend the inner product to a Hermitian one on $V_\CC$ and take $\theta(g) = (g^*)^{-1}$, where $g^*$ is the Hermitian adjoint of $g \in G(\CC)$. Then
$\theta(\bar g) = (g^T)^{-1}$ in a real orthonormal basis, so $G^{(\theta)} = G(\CC) \cap U(V_\CC)$, which is compact.
\end{proof}

\begin{remark}\label{algebraic-groups:remark-mumford-tate-reductive}
A variant of this argument, with the Weil operator of a polarisation in place of the transpose, shows that Mumford--Tate groups of polarisable Hodge structures are
reductive (Chapter~\ref{hodge-structures}).
\end{remark}

\section{Weil restriction of scalars}\label{algebraic-groups:section-weil-restriction}

\begin{definition}\label{algebraic-groups:definition-weil-restriction}
Let $L/k$ be a finite field extension and $Y$ an affine scheme over $L$. The \emph{Weil restriction} $\mathrm{Res}_{L/k}Y$ is the functor on $k$-algebras
$R \mapsto Y(R \otimes_kL)$.
\end{definition}

\begin{proposition}\label{algebraic-groups:proposition-weil-restriction}
$\mathrm{Res}_{L/k}Y$ is represented by an affine $k$-scheme. If $Y$ is a linear algebraic group of dimension $d$, then $\mathrm{Res}_{L/k}Y$ is a linear algebraic group of dimension
$d[L : k]$. If $L/k$ is Galois, then $(\mathrm{Res}_{L/k}Y)_L \cong \prod_{\sigma \in \mathrm{Gal}(L/k)}Y^\sigma$.
\end{proposition}

\begin{proof}
Choose a basis $e_1, \ldots, e_m$ of $L$ over $k$. For $Y = \Aa^n_L$, a point of $Y(R \otimes_kL)$ is an $n$-tuple $y_j = \sum_ix_{ij}e_i$ with $x_{ij} \in R$, so
$\mathrm{Res}_{L/k}\Aa^n_L = \Aa^{nm}_k$. If $Y = V(f_1, \ldots, f_s) \subseteq \Aa^n_L$, substitute $y_j = \sum_ix_{ij}e_i$ into $f_l$ and expand in the basis:
$f_l = \sum_iF_{li}(x)e_i$ with $F_{li} \in k[x]$. Then $\mathrm{Res}_{L/k}Y = V(F_{li}) \subseteq \Aa^{nm}_k$. The group structure is inherited functorially. For the dimension and the
Galois statement, use $(R \otimes_kL) \otimes_kL \cong \prod_\sigma R \otimes_kL$ for $R$ an $L$-algebra, which gives
$\mathrm{Res}_{L/k}Y(R) = \prod_\sigma Y^\sigma(R)$ for $L$-algebras $R$.
\end{proof}

\begin{example}\label{algebraic-groups:example-weil-restriction}
\begin{enumerate}
\item $\mathbb{S} = \mathrm{Res}_{\CC/\RR}\mathbb{G}_m$ is a two-dimensional torus over $\RR$ that is not split, with $\mathbb{S}_\CC \cong \mathbb{G}_m^2$; complex conjugation swaps the two
characters $z$ and $\bar z$.
\item For a number field $E$, $T_E = \mathrm{Res}_{E/\QQ}\mathbb{G}_m$ has $T_E(\QQ) = E^*$. If $E$ is a CM field, the Mumford--Tate group of an abelian variety with complex
multiplication by $E$ is a subtorus of $T_E$ (Section~\ref{algebraic-groups:section-abelian-varieties}).
\item The kernel $\mathbb{U}(1) = \ker(\mathrm{Nm} \colon \mathbb{S} \to \mathbb{G}_m)$ is the circle group, with $\mathbb{U}(1)(\RR) = S^1$, a torus that is compact and not split over $\RR$.
\end{enumerate}
\end{example}

\section{Exercises}\label{algebraic-groups:section-exercises}

\begin{exercise}\label{algebraic-groups:exercise-characters-gln}
Show that $X^*(\mathrm{GL}_n) = \ZZ$, generated by $\det$, and that $X^*(\mathrm{SL}_n) = 0$.
\end{exercise}

\begin{solution}
A character is a unit $\chi \in k[x_{ij}, 1/\det]$ with $\Delta(\chi) = \chi \otimes \chi$. The units of this ring are $c\det^m$, since $\det$ is irreducible, and group-likeness forces $c = 1$.
For $\mathrm{SL}_n$, the group $\mathrm{SL}_n(\bar k)$ is generated by the elementary matrices $1 + tE_{ij}$, $i \neq j$, and each family $t \mapsto 1 + tE_{ij}$ is a homomorphism
$\mathbb{G}_a \to \mathrm{SL}_n$. A character of $\mathbb{G}_a$ is a group-like unit of $k[t]$, hence $1$. So every character of $\mathrm{SL}_n$ is trivial on a generating set of
$\mathrm{SL}_n(\bar k)$, and as $\mathrm{SL}_n$ is reduced, it is trivial.
\end{solution}

\begin{exercise}\label{algebraic-groups:exercise-deligne-torus-examples}
Describe the representations of $\mathbb{S}$ on $\RR^2$ given by $h(a + ib) = \left(\begin{smallmatrix}a & -b \\ b & a\end{smallmatrix}\right)$ and by $h(z) = z\bar z \cdot \mathrm{id}$, in terms of
the decomposition of Proposition~\ref{algebraic-groups:proposition-deligne-torus}.
\end{exercise}

\begin{solution}
The first is multiplication by $z$ on $\CC = \RR^2$. On $V_\CC = \CC^2$ its eigenvectors are $(1, -i)$ with eigenvalue $z$ and $(1, i)$ with eigenvalue $\bar z$. With
$h(z)v = z^{-p}\bar z^{-q}v$, these are $V^{-1,0}$ and $V^{0,-1}$, a Hodge structure of weight $-1$ (the Hodge structure of $H_1$ of an elliptic curve). The second acts by
$z\bar z = z^1\bar z^1$ on all of $V_\CC$, so $V = V^{-1,-1}$, of weight $-2$.
\end{solution}

\begin{exercise}\label{algebraic-groups:exercise-unipotent-not-reductive}
Show directly that the group $B$ of invertible upper triangular $2 \times 2$ matrices has a representation that is not semisimple, and compute its unipotent radical.
\end{exercise}

\begin{solution}
The standard representation on $k^2$ has the invariant line $ke_1$, and an invariant complement would be spanned by a common eigenvector outside $ke_1$, which does not exist
because $\left(\begin{smallmatrix}1 & 1 \\ 0 & 1\end{smallmatrix}\right)$ has only the eigenline $ke_1$. The unipotent radical is the normal subgroup $U$ of unipotent upper triangular
matrices, and $B/U \cong \mathbb{G}_m^2$ is a torus.
\end{solution}

\section{Abelian Varieties}\label{algebraic-groups:section-abelian-varieties}
A complex abelian variety is a complex torus $\CC^g/\Lambda$ that embeds in projective space. Tori are the simplest compact K\"ahler manifolds, and their
cohomology is completely explicit: exterior powers of $H^1$. For this reason abelian varieties are the main testing ground of the Hodge conjecture. They
include the known cases (divisors, Lefschetz classes), the first instances of exceptional Hodge classes (Weil classes), and Deligne's theorem that Hodge classes
on abelian varieties are absolutely Hodge (Part~\ref{part-hodge-theory}). The following sections prove Riemann's criterion for a torus to be algebraic, describes line
bundles, polarisations, the dual and the Jacobian, and discusses endomorphisms and complex multiplication. References: \cite{birkenhake-lange},
\cite{mumford-abelian}, \cite[Chapter 2, Section 6]{griffiths-harris}.

\noindent\textbf{Prerequisites.} Chapter~\ref{curves} (Algebraic Curves and Riemann Surfaces), Chapter~\ref{divisors} (Divisors and Line Bundles),
Section~\ref{kahler:section-hodge-decomposition} (The Hodge Decomposition) and Chapter~\ref{algebraic-groups} (Linear Algebraic Groups).

Throughout, $V$ is a complex vector space of dimension $g$, $\Lambda \subseteq V$ a lattice (a discrete subgroup of rank $2g$ spanning $V$ over $\RR$), and $T = V/\Lambda$ the complex
torus, with its structure of compact complex Lie group (group law written additively). $\bar V^\vee$ denotes the space of conjugate-linear forms $V \to \CC$.

\section{Complex tori and their cohomology}\label{algebraic-groups:section-tori-abelian-varieties}

\begin{proposition}\label{algebraic-groups:proposition-torus-cohomology}
\begin{enumerate}
\item $H_1(T; \ZZ) = \Lambda$, and cup product gives ring isomorphisms $H^k(T; \ZZ) \cong \Lambda^k\Hom(\Lambda, \ZZ)$, the alternating $k$-forms on $\Lambda$.
\item The Hodge decomposition is $H^1(T; \CC) = \Hom_\RR(\Lambda \otimes \RR, \CC) = V^\vee \oplus \bar V^\vee$, with $H^{1,0} = V^\vee$ (spanned by the $dz_j$) and
$H^{0,1} = \bar V^\vee$ (spanned by the $d\bar z_j$), and $H^{p,q}(T) = \Lambda^pV^\vee \otimes \Lambda^q\bar V^\vee$.
\end{enumerate}
\end{proposition}

\begin{proof}
(1) $T$ is diffeomorphic to $(S^1)^{2g}$ via a $\ZZ$-basis of $\Lambda$, and $\RR^{2g} \to T$ is the universal covering, so $\pi_1(T) = H_1(T; \ZZ) = \Lambda$. The ring structure follows from
the K\"unneth formula (Theorem~\ref{singular-homology:theorem-kunneth}) and $H^\bullet(S^1) = \Lambda^\bullet\ZZ$. (2) With a flat metric the harmonic forms are the constant-coefficient
forms (Example~\ref{complex-manifolds:example-hodge-numbers}); a constant $1$-form is a real-linear form on $V$, and it splits uniquely into complex-linear and conjugate-linear parts.
\end{proof}

\begin{definition}\label{algebraic-groups:definition-abelian-variety}
A complex torus is an \emph{abelian variety} if it is projective. A homomorphism of tori is a holomorphic group homomorphism, and an \emph{isogeny} is a surjective homomorphism
with finite kernel. More generally, over any field an \emph{abelian variety} is a complete connected group variety; such a group is automatically commutative and projective
\cite[Chapter II]{mumford-abelian}.
\end{definition}

\begin{lemma}\label{algebraic-groups:lemma-homomorphisms}
Every holomorphic map $f \colon T \to T'$ of complex tori with $f(0) = 0$ is a homomorphism, induced by a complex-linear map $F \colon V \to V'$ with $F(\Lambda) \subseteq \Lambda'$.
\end{lemma}

\begin{proof}
Lift $f$ to $F \colon V \to V'$ with $F(0) = 0$. For $\lambda \in \Lambda$, $F(v + \lambda) - F(v) \in \Lambda'$ depends continuously on $v$, hence is constant. So the partial derivatives
of $F$ are $\Lambda$-periodic holomorphic functions, so they descend to holomorphic functions on the compact connected manifold $T$, which are constant by the maximum principle
(Proposition~\ref{several-complex-variables:proposition-identity-maximum}). So $F$ is linear, and $F(\Lambda) \subseteq \Lambda'$ because $F(\lambda) = F(0 + \lambda) - F(0) \in \Lambda'$.
\end{proof}

\begin{proposition}\label{algebraic-groups:proposition-torsion}
Let $T = V/\Lambda$ be a complex torus of dimension $g$ and $n \geq 1$.
\begin{enumerate}
\item Multiplication by $n$, $[n] \colon T \to T$, is a surjective homomorphism with kernel $T[n] = \frac1n\Lambda/\Lambda \cong (\ZZ/n\ZZ)^{2g}$. The torsion subgroup of $T$ is $(\Lambda \otimes \QQ)/\Lambda \cong (\QQ/\ZZ)^{2g}$, and it is
dense in $T$.
\item A homomorphism $f \colon T \to T'$ of tori of the same dimension is an isogeny if and only if its linear lift $F \colon V \to V'$ is an isomorphism; then $|\ker f| = [\Lambda' : F(\Lambda)]$.
\item If $f \colon T \to T'$ is an isogeny with $|\ker f| = d$, there is an isogeny $h \colon T' \to T$ with $h \circ f = [d]_T$ and $f \circ h = [d]_{T'}$. In particular being isogenous is an equivalence relation.
\end{enumerate}
\end{proposition}

\begin{proof}
(1) $[n]$ is induced by $v \mapsto nv$, which is bijective on $V$; so $[n]$ is surjective, and $x = v + \Lambda$ lies in the kernel if and only if $nv \in \Lambda$. A $\ZZ$-basis of $\Lambda$ identifies
$\frac1n\Lambda/\Lambda$ with $(\frac1n\ZZ/\ZZ)^{2g} \cong (\ZZ/n\ZZ)^{2g}$, and the union over $n$ is $(\Lambda \otimes \QQ)/\Lambda$. The rational points $\Lambda \otimes \QQ$ are dense in $\Lambda \otimes \RR = V$, so
their image is dense in $T$.

(2) By Lemma~\ref{algebraic-groups:lemma-homomorphisms}, $f$ is induced by a $\CC$-linear $F$ with $F(\Lambda) \subseteq \Lambda'$. If $F$ is an isomorphism, $F(\Lambda)$ is a lattice in $V'$ contained in $\Lambda'$, hence of
finite index, $f$ is surjective, and $\ker f = F^{-1}(\Lambda')/\Lambda \cong \Lambda'/F(\Lambda)$ is finite. Conversely, if $f$ has finite kernel, then $\ker F$ is a complex subspace whose image in $T$ is contained in the
finite group $\ker f$; a connected subspace mapping into a finite set maps to a point, so $\ker F \subseteq \Lambda$, which forces $\ker F = 0$; equal dimensions make $F$ an isomorphism.

(3) The group $\Lambda'/F(\Lambda)$ has order $d$, so $d\Lambda' \subseteq F(\Lambda)$. Hence $H = dF^{-1} \colon V' \to V$ maps $\Lambda'$ into $\Lambda$ and induces a homomorphism $h \colon T' \to T$ with
$HF = d$ and $FH = d$, which is an isogeny by (2). Reflexivity and transitivity of isogeny are clear, and (3) gives symmetry.
\end{proof}

\section{Line bundles and the Riemann conditions}\label{algebraic-groups:section-riemann}

\begin{proposition}\label{algebraic-groups:proposition-neron-severi}
Identify $H^2(T; \ZZ)$ with the alternating forms $E \colon \Lambda \times \Lambda \to \ZZ$ (Proposition~\ref{algebraic-groups:proposition-torus-cohomology}). A class $E$ lies in
$H^{1,1}(T)$ if and only if $E(iu, iv) = E(u, v)$ for all $u, v \in V$ (extending $E$ $\RR$-bilinearly), if and only if $E = \operatorname{Im}H$ for a (unique) Hermitian form $H$ on $V$,
namely $H(u, v) = E(iu, v) + iE(u, v)$. Every such $E$ is $c_1(L)$ for a holomorphic line bundle $L$, so
\[
\mathrm{NS}(T) := \operatorname{im}(c_1) = \{H \text{ Hermitian on } V : \operatorname{Im}H(\Lambda, \Lambda) \subseteq \ZZ\}.
\]
\end{proposition}

\begin{proof}
A real alternating form $E$ on $V$ is a constant real $2$-form. Its $(2,0)$ and $(0,2)$ components vanish exactly when $E$ is of type $(1,1)$ for the complex structure $i$,
that is, $E(iu, iv) = E(u, v)$: the $(2,0) + (0,2)$ part is $\frac12(E - E(i\cdot, i\cdot))$. Given this, $H(u, v) = E(iu, v) + iE(u, v)$ is Hermitian:
$H(iu, v) = E(-u, v) + iE(iu, v) = i(E(iu, v) + iE(u, v)) = iH(u, v)$, and $\overline{H(v, u)} = E(iv, u) - iE(v, u) = E(iu, v) + iE(u, v)$, using $E(iv, u) = -E(u, iv) = -E(iu, i^2v) = E(iu, v)$.
Conversely $\operatorname{Im}$ of a Hermitian form satisfies the condition. The last statement is Lemma~\ref{kodaira:lemma-integral-11}.
\end{proof}

\begin{theorem}[Riemann]\label{algebraic-groups:theorem-riemann-conditions}
A complex torus $T = V/\Lambda$ is an abelian variety if and only if there is a positive definite Hermitian form $H$ on $V$ with $\operatorname{Im}H(\Lambda, \Lambda) \subseteq \ZZ$.
Such an $H$ (or its alternating form $E = \operatorname{Im}H$) is a \emph{polarisation}; it is the first Chern class of an ample line bundle.
\end{theorem}

\begin{proof}
For a constant real $2$-form $\omega$ put $E_\omega = -\omega$ and $H_\omega(u, v) = E_\omega(iu, v) + iE_\omega(u, v)$. If $\omega$ is of type $(1,1)$, then $H_\omega$ is Hermitian
(Proposition~\ref{algebraic-groups:proposition-neron-severi}) with $\operatorname{Re}H_\omega(u, u) = -\omega(iu, u) = \omega(u, iu)$. So $\omega$ is a positive $(1,1)$-form, that is, the
fundamental form of a Hermitian metric (Proposition~\ref{kahler:proposition-fundamental-form-coordinates}), if and only if $H_\omega$ is positive definite. Conversely every Hermitian $H$ is
$H_\omega$ for $\omega = -\operatorname{Im}H$.

Suppose $H$ is positive definite with $\operatorname{Im}H$ integral on $\Lambda$. Then $\omega = -\operatorname{Im}H$ is a positive constant $(1,1)$-form, closed because its coefficients
are constant, and its class is integral. So $T$ is a Hodge manifold (Definition~\ref{kodaira:definition-hodge-manifold}) and projective by
Corollary~\ref{kodaira:corollary-projective-criterion}.

Conversely, let $T$ be projective, with an ample line bundle $L$ and a metric of positive curvature form $\omega_h$ (Theorem~\ref{gaga:theorem-projective-kahler}). Average over translations:
$\bar\omega = \int_Tt_x^*\omega_h\,dx$ for the normalised Haar measure (Theorem~\ref{lie-groups:theorem-haar}). Each $t_x^*\omega_h$ is positive and cohomologous to $\omega_h$, because
translations are homotopic to the identity, so $\bar\omega$ is a positive translation-invariant closed $(1,1)$-form in the integral class $c_1(L)$. A translation-invariant form has constant
coefficients. So $H_{\bar\omega}$ is positive definite, and $\operatorname{Im}H_{\bar\omega} = -\bar\omega$ is integral on $\Lambda$, being (up to sign) the integral class $c_1(L)$ in
$H^2(T; \ZZ) = \Lambda^2\Hom(\Lambda, \ZZ)$.
\end{proof}

\begin{remark}\label{algebraic-groups:remark-sign-convention}
The proof shows that ample line bundles correspond to positive definite $H$ once a constant form $\omega$ is matched with the alternating form $-\omega$; this is the convention of
\cite{birkenhake-lange}, where $c_1(L) = \operatorname{Im}H$. With the opposite matching one gets negative definite forms instead; Riemann's condition is insensitive to this, since
$-H$ is positive definite if $H$ is negative definite.
\end{remark}

\begin{proposition}\label{algebraic-groups:proposition-period-matrix}
Let $E$ be a polarisation. There is a basis $\lambda_1, \ldots, \lambda_g, \mu_1, \ldots, \mu_g$ of $\Lambda$ with
$E(\lambda_j, \mu_k) = d_j\delta_{jk}$ and $E(\lambda_j, \lambda_k) = E(\mu_j, \mu_k) = 0$, where $d_1 \mid d_2 \mid \cdots \mid d_g$ are positive integers, the \emph{type} of the polarisation.
It is \emph{principal} if all $d_j = 1$. Then $\frac{1}{d_1}\mu_1, \ldots, \frac{1}{d_g}\mu_g$ is a $\CC$-basis of $V$, and in it $\Lambda = Z\ZZ^g + D\ZZ^g$ with $D = \mathrm{diag}(d_j)$
and $Z$ a symmetric complex $g \times g$ matrix with $\operatorname{Im}Z$ positive definite. Conversely every such $Z$ in the \emph{Siegel upper half space} $\mathfrak{H}_g$ defines a
polarised abelian variety of type $D$.
\end{proposition}

\begin{proof}
The symplectic basis is the elementary divisor theorem for alternating forms over $\ZZ$ (Frobenius), proved like the Smith normal form (Theorem~\ref{rings:theorem-smith}); see
\cite{birkenhake-lange}. Since $H$ is positive definite, the real span of the $\mu_j$ is $E$-isotropic, and $H$ is real and positive on it, so the $\mu_j$ are $\CC$-linearly
independent: a complex relation $\sum c_j\mu_j = 0$ with $c = a + ib$ gives $\sum a_j\mu_j = -i\sum b_j\mu_j$, and pairing with $H$ forces $a = b = 0$. Writing
$\lambda_j = \sum_kZ_{kj}\mu_k/d_k$, the conditions that the $\lambda$'s span an isotropic subspace and that $H$ is positive definite translate into $Z = Z^T$ and
$\operatorname{Im}Z > 0$ by a direct computation \cite{birkenhake-lange}. The converse is the same computation read backwards together with
Theorem~\ref{algebraic-groups:theorem-riemann-conditions}.
\end{proof}

\begin{example}\label{algebraic-groups:example-dimension-one}
For $g = 1$, $\mathfrak{H}_1$ is the upper half plane, and $\Lambda = \ZZ\tau + \ZZ$ with $\operatorname{Im}\tau > 0$. Every one-dimensional torus is an abelian variety, as Riemann's
condition only asks $\operatorname{Im}\tau > 0$ (compare Example~\ref{kahler:example-curves-riemann}). The alternating form $E(\tau, 1) = 1$ gives
$H(u, v) = u\bar v/\operatorname{Im}\tau$.
\end{example}

\begin{counterexample}\label{algebraic-groups:counterexample-non-algebraic-torus}
For $g \geq 2$ the general complex torus is not algebraic. Fix a real basis identifying $V$ with $\RR^{2g}$ and $\Lambda$ with $\ZZ^{2g}$, so that complex tori correspond to complex
structures $J$ on $\RR^{2g}$. For each of the countably many nonzero integral alternating forms $E$, the condition $E(Ju, Jv) = E(u, v)$ defines a proper closed real-analytic subset of the
space of complex structures when $g \geq 2$ \cite{birkenhake-lange}. Outside the countable union of these subsets,
$\mathrm{NS}(T) = 0$ by Proposition~\ref{algebraic-groups:proposition-neron-severi}, so there is no ample line bundle and $T$ is not projective. For such tori the Hodge conjecture
in degree $2$ has no analytic cycles to work with; more generally the analytic version of the conjecture fails for some complex tori (Chapter~\ref{known-cases}).
\end{counterexample}

\begin{theorem}[Appell--Humbert]\label{algebraic-groups:theorem-appell-humbert}
Holomorphic line bundles on $T$ correspond bijectively to pairs $(H, \alpha)$ with $H$ Hermitian, $E = \operatorname{Im}H$ integral on $\Lambda$, and $\alpha \colon \Lambda \to U(1)$ a
\emph{semicharacter}: $\alpha(\lambda + \mu) = \alpha(\lambda)\alpha(\mu)e^{i\pi E(\lambda, \mu)}$. The bundle $L(H, \alpha)$ is the quotient of $V \times \CC$ by
$\lambda \cdot (v, t) = (v + \lambda, \alpha(\lambda)e^{\pi H(v, \lambda) + \frac{\pi}{2}H(\lambda, \lambda)}t)$, and $c_1(L(H, \alpha)) = E$. Its sections are the
\emph{theta functions}, and if $H$ is positive definite of type $(d_1, \ldots, d_g)$, then $h^0(L) = d_1 \cdots d_g$. Moreover $L^2$ is globally generated and $L^3$ is very ample
(Lefschetz).
\end{theorem}

This is \cite{birkenhake-lange}; see also \cite{mumford-abelian}.

\section{The dual abelian variety}\label{algebraic-groups:section-dual}

\begin{proposition}\label{algebraic-groups:proposition-dual}
$\Pic^0(T) \cong \hat T = \bar V^\vee/\hat\Lambda$, where $\hat\Lambda = \{l \in \bar V^\vee : \operatorname{Im}l(\Lambda) \subseteq \ZZ\}$, a complex torus of dimension $g$, the
\emph{dual torus}. A polarisation $H$ defines a homomorphism $\varphi_H \colon T \to \hat T$, $v \mapsto H(v, \cdot)$, which is an isogeny of degree $(d_1 \cdots d_g)^2$, and an
isomorphism if the polarisation is principal. If $T$ is an abelian variety, so is $\hat T$.
\end{proposition}

\begin{proof}
By the exponential sequence, $\Pic^0(T) = H^1(T, \mathcal{O})/H^1(T; \ZZ)$ (Theorem~\ref{complex-manifolds:theorem-exponential}), and $H^1(T, \mathcal{O}) = H^{0,1} = \bar V^\vee$
(Proposition~\ref{algebraic-groups:proposition-torus-cohomology}). The map $H^1(T; \RR) = \Hom_\RR(V, \RR) \to \bar V^\vee$ sends a real linear form $a$ to its conjugate-linear part
$l_a$, and $a = 2\operatorname{Re}l_a$. So the image of $H^1(T; \ZZ) = \Hom(\Lambda, \ZZ)$ is $\{l : 2\operatorname{Re}l(\Lambda) \subseteq \ZZ\}$, which the $\CC$-linear automorphism
$l \mapsto 2il$ of $\bar V^\vee$ carries onto $\hat\Lambda$, since $\operatorname{Im}(2il) = 2\operatorname{Re}l$. The map
$v \mapsto H(v, \cdot)$ is complex-linear from $V$ to $\bar V^\vee$ and maps $\Lambda$ into $\hat\Lambda$ because $\operatorname{Im}H(\Lambda, \Lambda) \subseteq \ZZ$. It is injective since
$H$ is nondegenerate, so its kernel on $T$ is $\hat\Lambda'/\Lambda$ with $\hat\Lambda' = \{v : \operatorname{Im}H(v, \Lambda) \subseteq \ZZ\}$, of order $\det E = (d_1 \cdots d_g)^2$ in
a symplectic basis. The inverse of $H$ gives a polarisation on $\hat T$ \cite{birkenhake-lange}.
\end{proof}

\section{Jacobians}\label{algebraic-groups:section-jacobians}

\begin{theorem}\label{algebraic-groups:theorem-jacobian}
Let $C$ be a compact Riemann surface of genus $g$. The Jacobian $J(C) = H^0(C, \Omega^1)^\vee/H_1(C; \ZZ)$ (Definition~\ref{curves:definition-jacobian}) is a principally
polarised abelian variety, the polarisation being the intersection form on $H_1(C; \ZZ)$.
\end{theorem}

\begin{proof}
$J(C)$ is a complex torus of dimension $g$ (Proposition~\ref{curves:proposition-jacobian-torus}). The intersection form $E$ on $H_1(C; \ZZ)$ is unimodular (Poincar\'e duality,
Corollary~\ref{poincare-duality:corollary-cup-pairing}), so a symplectic basis $a_1, \ldots, a_g, b_1, \ldots, b_g$ exists with $E(a_j, b_k) = \delta_{jk}$. Extend $E$ $\RR$-linearly to
$V = H^0(\Omega^1)^\vee \cong H_1(C; \RR)$. Under Poincar\'e duality $H_1(C; \RR) \cong H^1(C; \RR)$ and the real isomorphism $H^1(C; \RR) \to H^{1,0}$, $\alpha \mapsto \alpha^{1,0}$,
the form $E$ becomes $\pm\int_C\alpha \wedge \beta = \pm 2\operatorname{Re}\int_C\alpha^{1,0} \wedge \overline{\beta^{1,0}}$, and the complex structure of $V$ becomes that of $H^{1,0}$ (up to
conjugation). So $E(iu, iv) = E(u, v)$, and the Hermitian form $H(u, v) = E(iu, v) + iE(u, v)$ is, up to sign and a positive factor, $i\int_C\alpha^{1,0} \wedge \overline{\beta^{1,0}}$, which
is definite by the Hodge--Riemann relations (Theorem~\ref{kahler:theorem-hodge-riemann}, Example~\ref{kahler:example-curves-riemann}). Replacing $E$ by $-E$ if
necessary, $H$ is positive definite, and $J(C)$ is an abelian variety by Theorem~\ref{algebraic-groups:theorem-riemann-conditions}. The polarisation is principal as $E$ is unimodular.
\end{proof}

\begin{remark}\label{algebraic-groups:remark-torelli-schottky}
By Torelli's theorem the principally polarised Jacobian determines $C$. The Jacobians form a subvariety of dimension $3g - 3$ of the moduli space $\mathfrak{H}_g/\mathrm{Sp}_{2g}(\ZZ)$ of
principally polarised abelian varieties, which has dimension $g(g + 1)/2$. For $g \leq 3$ the dimensions agree and the general principally polarised abelian variety is a Jacobian;
for $g \geq 4$ it is not, and characterising Jacobians is the \emph{Schottky problem} \cite{birkenhake-lange}.
\end{remark}

\section{Endomorphisms and complex multiplication}\label{algebraic-groups:section-endomorphisms}

\begin{proposition}\label{algebraic-groups:proposition-elliptic-endomorphisms}
Let $E = \CC/(\ZZ\tau + \ZZ)$. Then $\End(E)$ is $\ZZ$, or an order in the imaginary quadratic field $\QQ(\tau)$; the latter happens if and only if $\tau$ is quadratic over $\QQ$.
\end{proposition}

\begin{proof}
By Lemma~\ref{algebraic-groups:lemma-homomorphisms} an endomorphism is multiplication by some $\alpha \in \CC$ with $\alpha\Lambda \subseteq \Lambda$. Then $\alpha = a + b\tau$ and
$\alpha\tau = c + d\tau$ with integers $a, b, c, d$. If $b \neq 0$, eliminating $\alpha$ gives $b\tau^2 + (a - d)\tau - c = 0$, so $\tau$ is quadratic, imaginary as $\operatorname{Im}\tau > 0$,
and $\alpha \in \QQ(\tau)$ is integral (it is an eigenvalue of an integer matrix). If $b = 0$ then $\alpha = a \in \ZZ$. Conversely, if $a\tau^2 + b\tau + c = 0$ with integers $a \neq 0$,
$b$, $c$, then $\alpha = a\tau$ satisfies $\alpha \cdot 1 = a\tau \in \Lambda$ and $\alpha\tau = -b\tau - c \in \Lambda$, so multiplication by $\alpha \notin \ZZ$ is an endomorphism.
\end{proof}

\begin{definition}\label{algebraic-groups:definition-cm}
An abelian variety $A$ of dimension $g$ has \emph{complex multiplication} (CM) if $\End(A) \otimes \QQ$ contains a commutative semisimple subalgebra of dimension $2g$ over $\QQ$,
for instance a CM field $F$ (a totally imaginary quadratic extension of a totally real field) with $[F : \QQ] = 2g$.
\end{definition}

\begin{theorem}\label{algebraic-groups:theorem-endomorphism-algebra}
Let $A$ be an abelian variety over $\CC$.
\begin{enumerate}
\item (Poincar\'e reducibility) $A$ is isogenous to a product $A_1^{n_1} \times \cdots \times A_r^{n_r}$ with the $A_i$ simple and pairwise non-isogenous, and
$\End(A) \otimes \QQ \cong \prod_iM_{n_i}(D_i)$ with $D_i = \End(A_i) \otimes \QQ$ division algebras.
\item A polarisation defines the \emph{Rosati involution} $f \mapsto f^\dagger$ on $\End(A) \otimes \QQ$, and $\operatorname{tr}(ff^\dagger) > 0$ for $f \neq 0$.
\item (Shimura--Taniyama) If $A$ has CM by a CM field $F$ of degree $2g$, then $A \cong \CC^g/\Phi(\mathfrak{a})$ for a fractional ideal $\mathfrak{a}$ of $F$ and a \emph{CM type}
$\Phi = \{\phi_1, \ldots, \phi_g\}$, a set of embeddings $F \to \CC$ containing exactly one of each pair of complex conjugate embeddings.
\end{enumerate}
\end{theorem}

This is \cite{birkenhake-lange} and \cite{mumford-abelian}. The endomorphism algebras $D_i$ are classified by
Albert into four types. For the Hodge conjecture, CM abelian varieties are the case where the Mumford--Tate group is a torus (Chapter~\ref{hodge-structures}), and where exceptional
Hodge classes first appear (Chapter~\ref{known-cases}).

\begin{example}\label{algebraic-groups:example-cm}
\begin{enumerate}
\item $E_i = \CC/\ZZ[i]$ has $\End(E_i) = \ZZ[i]$ and $E_\rho = \CC/\ZZ[\rho]$, $\rho = e^{2\pi i/3}$, has $\End(E_\rho) = \ZZ[\rho]$.
\item For $E$ without CM, $\End(E \times E) \otimes \QQ = M_2(\QQ)$, and $\mathrm{NS}(E \times E) \otimes \QQ$ has rank $3$; for $E$ with CM it has rank $4$. This is
Example~\ref{divisors:example-picard-numbers}(3).
\item The Jacobian of the curve $y^2 = x^5 - 1$ has CM by $\QQ(\zeta_5)$, a CM field of degree $4 = 2g$.
\end{enumerate}
\end{example}

\section{Exercises on Abelian Varieties}\label{algebraic-groups:section-exercises-abelian-varieties}

\begin{exercise}\label{algebraic-groups:exercise-ns-product}
Let $E_1$, $E_2$ be elliptic curves. Show that $\mathrm{NS}(E_1 \times E_2) \cong \ZZ^2 \oplus \Hom(E_1, E_2)$, and deduce that its rank is $2$, $3$ or $4$.
\end{exercise}

\begin{solution}
Write $\Lambda = \Lambda_1 \oplus \Lambda_2$. An integral alternating form of type $(1,1)$ on $\Lambda$ restricts to such forms on $\Lambda_1$ and $\Lambda_2$, each a multiple of the generator
($\ZZ^2$), and its mixed part $E(\lambda_1, \lambda_2)$ is a pairing of type $(1,1)$ between $\Lambda_1$ and $\Lambda_2$, which via $E_1$'s principal polarisation is the same as a
$\ZZ$-linear map $\Lambda_1 \to \Lambda_2$ commuting with $i$ after identification, that is, a homomorphism $E_1 \to E_2$ (Lemma~\ref{algebraic-groups:lemma-homomorphisms}). By
Proposition~\ref{algebraic-groups:proposition-elliptic-endomorphisms} and its analogue for homomorphisms, $\Hom(E_1, E_2)$ has rank $0$, $1$ or $2$.
\end{solution}

\begin{exercise}\label{algebraic-groups:exercise-hodge-classes-torus}
Show that the Hodge classes in $H^2(T; \QQ)$ of a complex torus form the space $\mathrm{NS}(T) \otimes \QQ$, and that for $T = E^g$, $E$ without CM, this space has dimension
$g(g + 1)/2$.
\end{exercise}

\begin{solution}
The first statement is Proposition~\ref{algebraic-groups:proposition-neron-severi}. For $E^g$, $\mathrm{NS} \otimes \QQ$ is the space of $g \times g$ Hermitian matrices with entries in
$\End(E) \otimes \QQ = \QQ$ that are Hermitian for the Rosati involution, that is, symmetric rational matrices, of dimension $g(g + 1)/2$ \cite{birkenhake-lange}.
\end{solution}

\begin{exercise}\label{algebraic-groups:exercise-isogeny-dual}
Show that multiplication by $n$ on $T$ is an isogeny of degree $n^{2g}$ with kernel $\frac1n\Lambda/\Lambda \cong (\ZZ/n)^{2g}$.
\end{exercise}

\begin{solution}
The map is induced by $v \mapsto nv$ on $V$, surjective, with kernel $\{v : nv \in \Lambda\}/\Lambda = \frac1n\Lambda/\Lambda$, which has $n^{2g}$ elements since $\Lambda \cong \ZZ^{2g}$.
\end{solution}
